%--------------------------------------------------------------------------------
%
%
%--------------------------------------------------------------------------------
\pagebreak
\section*{ADD — Integer Addition}
\addcontentsline{toc}{section}{ADD - Integer Addition}

\instructionentry{Syntax}{
  \texttt{ADD RegR, RegB, RegA} \\
  \texttt{ADD RegR, RegB, Immed15} \\
  \texttt{ADD RegR, Immed13( RegB )} \\
  \texttt{ADD RegR, RegX ( RegB )}
}

\instructionentry{Format}{

	The ADD instruction uses the register, the immediate, the indexed and the register indexed instruction formats. \\

   	\begin{flushleft}
    	\begin{tikzpicture}[scale=0.7, transform shape]
        
        \draw[help lines, gray!70, dashed] (0,0) grid(16,1.5);
        
        \node[  tsRectangle, 
                minimum width=16cm,
                minimum height=1cm,
                text width=3cm,
                text centered,
                fill=white!50] (blocks) at (8,1) { };
                
     	\draw[-] (1,0.5) -- (1, 1.5);
     	\draw[-] (3,0.5) -- (3, 1.5);
     	\draw[-] (5,0.5) -- (5, 1.5);
     	\draw[-] (6.5,0.5) -- (6.5, 1.5);
     	\draw[-] (8.5,0.5) -- (8.5, 1.5);
     	\draw[-] (9.5,0.5) -- (9.5, 1.5);
     	\draw[-] (11.5,0.5) -- (11.5, 1.5);
     	
     	\node at (0.5,0.25) {2};
    	\node at (2,0.25) {4};
    	\node at (4,0.25) {4};
    	\node at (5.75,0.25) {3};
    	\node at (7.5,0.25) {4};
    	\node at (9,0.25) {2};
    	\node at (10.5,0.25) {4};
   	 	\node at (13.75,0.25) {9};
   	 	
   	 	\node at (0.5,1) {\small 0};
		\node at (2,1) {\small 1};
		\node at (4,1) {\small Reg R};
		\node at (5.75,1) {\small 0};
		\node at (7.5,1) {\small Reg B};
		\node at (9,1) {\small Opt 2};
		\node at (10.5,1) {\small Reg A};
		\node at (13.75,1) {\small S-IMM-9 / special };
   	 	
       \end{tikzpicture}
	\end{flushleft}
 	\medskip
 	\begin{flushleft}
    	\begin{tikzpicture}[scale=0.7, transform shape]
        
        \draw[help lines, gray!70, dashed] (0,0) grid(16,1.5);
        
        \node[  tsRectangle, 
                minimum width=16cm,
                minimum height=1cm,
                text width=3cm,
                text centered,
                fill=white!50] (blocks) at (8,1) { };
                
     	\draw[-] (1,0.5) -- (1, 1.5);
     	\draw[-] (3,0.5) -- (3, 1.5);
     	\draw[-] (5,0.5) -- (5, 1.5);
     	\draw[-] (6.5,0.5) -- (6.5, 1.5);
     	\draw[-] (8.5,0.5) -- (8.5, 1.5);
          	
     	\node at (0.5,0.25) {2};
    	\node at (2,0.25) {4};
    	\node at (4,0.25) {4};
    	\node at (5.75,0.25) {3};
    	\node at (7.5,0.25) {4};
   	 	\node at (11.5,0.25) {15};
   	 	
   	 	\node at (0.5,1) {\small 0};
		\node at (2,1) {\small 1};
		\node at (4,1) {\small Reg R};
		\node at (5.75,1) {\small 0};
		\node at (7.5,1) {\small Reg B};
		\node at (11.5,1) {\small S-IMM-15};
   	 	
       \end{tikzpicture}
	\end{flushleft}
	\medskip
	\begin{flushleft}
    \begin{tikzpicture}[scale=0.7, transform shape]
        
        \draw[help lines, gray!70, dashed] (0,0) grid(16,1.5);
        
        \node[  tsRectangle, 
                minimum width=16cm,
                minimum height=1cm,
                text width=3cm,
                text centered,
                fill=white!50] (blocks) at (8,1) { };
                
     	\draw[-] (1,0.5) -- (1, 1.5);
     	\draw[-] (3,0.5) -- (3, 1.5);
     	\draw[-] (5,0.5) -- (5, 1.5);
     	\draw[-] (6.5,0.5) -- (6.5, 1.5);
     	\draw[-] (8.5,0.5) -- (8.5, 1.5);
     	\draw[-] (9.5,0.5) -- (9.5, 1.5);
     	
     	\node at (0.5,0.25) {2};
    	\node at (2,0.25) {4};
    	\node at (4,0.25) {4};
    	\node at (5.75,0.25) {3};
    	\node at (7.5,0.25) {4};
    	\node at (9,0.25) {2};
     	\node at (12.25,0.25) {13};
   	 	
   	 	\node at (0.5,1) {\small 1};
		\node at (2,1) {\small 1};
		\node at (4,1) {\small Reg R};
		\node at (5.75,1) {\small 0};
		\node at (7.5,1) {\small Reg B};
		\node at (9,1) {\small dw};
		\node at (12.25,1) {\small S-IMM-13 / special};
   	 	
       \end{tikzpicture}
	\end{flushleft}
 	\medskip
	\begin{flushleft}
    \begin{tikzpicture}[scale=0.7, transform shape]
        
        \draw[help lines, gray!70, dashed] (0,0) grid(16,1.5);
        
        \node[  tsRectangle, 
                minimum width=16cm,
                minimum height=1cm,
                text width=3cm,
                text centered,
                fill=white!50] (blocks) at (8,1) { };
                
     	\draw[-] (1,0.5) -- (1, 1.5);
     	\draw[-] (3,0.5) -- (3, 1.5);
     	\draw[-] (5,0.5) -- (5, 1.5);
     	\draw[-] (6.5,0.5) -- (6.5, 1.5);
     	\draw[-] (8.5,0.5) -- (8.5, 1.5);
     	\draw[-] (9.5,0.5) -- (9.5, 1.5);
     	\draw[-] (11.5,0.5) -- (11.5, 1.5);
     	
     	\node at (0.5,0.25) {2};
    	\node at (2,0.25) {4};
    	\node at (4,0.25) {4};
    	\node at (5.75,0.25) {3};
    	\node at (7.5,0.25) {4};
    	\node at (9,0.25) {2};
    	\node at (10.5,0.25) {4};
   	 	\node at (13.75,0.25) {9};
   	 	
   	 	\node at (0.5,1) {\small 1};
		\node at (2,1) {\small 1};
		\node at (4,1) {\small Reg R};
		\node at (5.75,1) {\small 0};
		\node at (7.5,1) {\small Reg B};
		\node at (9,1) {\small dw};
		\node at (10.5,1) {\small Reg A};
		\node at (13.75,1) {\small 0};
   	 	
       \end{tikzpicture}
	\end{flushleft}
}

\instructionentry{Description}{

The \texttt{ADD} instruction adds two operands and stores them to the target register \texttt{RegR}. The immediate instruction format will add the sign extended 15-bit value to \texttt{RegB}. The register instruction format instruction will add \texttt{RegA} and \texttt{RegB}. 

The indexed instruction format will load the content of memory address formed by adding the scaled immediate 13-bit field to the base register \texttt{RegB} and add the data value to \texttt{RegA}. The register indexed instruction format will load the content of memory address formed by adding \texttt{RegX} to the base register \texttt{RegB} and add the data value to \texttt{RegA}. Both load formats will use the dw field to specify the data value width. A dw value of zero load an 8-bot, a value of 1 a 16-bit value,  value of 2 a 32-bit value and a value of 3 a 64-bit value. The data value fetched is signed extended to a 64-bit value.

}

\instructionentry{Operation}{
  	\texttt{RegR <- RegB + RegA} {(register format)}\\
  	\smallskip
  	\texttt{RegR <- RegB + immOperand( )} {(immediate format)}\\
   	\smallskip
  	\texttt{RegR <- RegR + memOperand( )} {(indexed formats)}\\
}

\instructionentry{Exceptions}{

 	Integer Overflow \\
	Data Alignment \\
	Memory Access Violation \\
	Memory Protection Violation \\
	Data TLB miss
}

\instructionentry{Notes}{

	None. 
}




